\clearpage
\begin{summarybox}
	\textbf{\textcolor{CSummary}{一句话核心}}

	\textbf{函数图像对称变换的实质是点的坐标替换，对应解析式按对称规则改变自变量或因变量的符号或位置。}

	\medskip
	\textbf{\textcolor{CSummary}{使用条件}}
	\begin{itemize}[leftmargin=*, itemsep=0.5em, topsep=0.4em]
		\item \textbf{条件 1（原函数存在）：} 存在已知函数 $y=f(x)$ 及其图像。
		\item \textbf{条件 2（对称形式确定）：} 明确对称轴或对称中心（包括 $x$ 轴、$y$ 轴、原点、直线 $y=x$）。
	\end{itemize}

	\medskip
	\textbf{\textcolor{CSummary}{关键提醒}}
	\begin{itemize}[leftmargin=*, itemsep=0.5em, topsep=0.4em]
		\item \textbf{易错点（符号混淆）：} 注意区分关于 $x$ 轴和 $y$ 轴对称的公式，符号位置不同。
		\item \textbf{检查项（反函数存在）：} 关于直线 $y=x$ 对称时，必须验证原函数可逆，否则不能直接写反函数。
	\end{itemize}
\end{summarybox}
\clearpage
\begin{summarybox}
\textbf{补充速查｜先认两幅图，再认对称位置}

\[\begin{array}{c|c}
\text{对称位置}&\text{由 }f\text{ 得到的新函数 }g\\\hline
x=a & g(x)=f(2a-x)\\
y=b & g(x)=2b-f(x)\\
(a,b) & g(x)=2b-f(2a-x)
\end{array}\]
\[\boxed{y=f(a+x)\text{ 与 }y=f(b-x)\text{ 关于 }x=\frac{b-a}{2}\text{ 相互对称。}}\]
\textbf{三个最后检查：}定义域是否完整对应？反函数是否存在？两幅图相互对称，还是同一幅图反射后与自身重合？

\textbf{最快回验：}猜轴 $x=c$ 后，检查 $F(2c-x)=G(x)$ 和定义域；猜中心 $(a,b)$ 后，检查 $G(x)=2b-F(2a-x)$ 和定义域。这些条件都须覆盖整幅图像。

\end{summarybox}
